\documentclass[11pt]{article}
\usepackage[margin=25mm]{geometry}
\usepackage{fontspec}
\setmainfont{Latin Modern Roman}
\usepackage{amsmath,amssymb,mathtools}
\usepackage{unicode-math}
\setmathfont{Latin Modern Math}
\usepackage{xurl,hyperref}
\hypersetup{colorlinks=true,urlcolor=blue}
\setcounter{secnumdepth}{0}
\setlength{\parindent}{0pt}
\setlength{\parskip}{5pt}
\setlength{\emergencystretch}{3em}
\newcommand{\sref}[2]{\hyperlink{source:#1:#2}{[S#2]}}
\newcommand{\cataloguescope}[1]{\paragraph{Recorded target.}#1}
\newcommand{\R}{\mathbb R}
\newcommand{\C}{\mathbb C}
\newcommand{\N}{\mathbb N}
\newcommand{\Z}{\mathbb Z}
\newcommand{\Q}{\mathbb Q}
\newcommand{\F}{\mathbb F}
\newcommand{\eps}{\varepsilon}
\begin{document}
\section{Geometric Langlands and Quantum Physics}\label{statement}

\subsection{Definitions and mathematical statement}

\paragraph{Shared notation.}$\mathbb N=\{1,2,\ldots\}$; the standard number fields have their usual meanings. All additional conventions are stated locally.


\paragraph{Definitions and notation.}The geometric Langlands programme relates geometric objects associated with a reductive group and its Langlands dual, such as local systems and sheaf-theoretic data on moduli spaces of bundles over a curve. In quantum field theory a duality relates descriptions of physical observables or state spaces. The challenge is to make connections between such mathematical correspondences and fundamental physical symmetries precise.

\paragraph{Catalogue wording retained for traceability.}
Connect the Langlands program to fundamental physics symmetries.
\paragraph{Scope and source qualification.}This broad research objective fixes neither a group, space-time theory nor single proposed equivalence. The historical DARPA mirror was unavailable during this check.
\subsection{Short English statement}
Explain rigorous connections between Langlands correspondences and symmetries or dualities in fundamental physics.
\subsection{Research attempt}
\paragraph{Model, date and investigated scope.}
OpenAI GPT-6 Astra (configured model \texttt{gpt-6-astra}), ultra reasoning; generated 11 October 2026. Verified lattice and quadratic-form compatibility of rank-one electric-magnetic exchange, with a torus lattice formulation. No quantum or geometric Langlands equivalence is constructed.

\paragraph{Sources and literature check.}
Read the supplied complete definition, catalogue wording, scope qualifications and source list for this canonical ID. The original-source citations are inherited as supplied references; no fresh primary-source or novelty search was performed for this attempt. No current open-status claim is inferred from this calculation.

\paragraph{An abelian lattice calculation.}A modest place to test a symmetry interpretation is a torus \(T=(S^1)^r\). Its character lattice \(X^*(T)\) consists of continuous homomorphisms \(T\to S^1\), and its cocharacter lattice \(X_*(T)\) consists of homomorphisms \(S^1\to T\); both identify with \(\mathbb Z^r\). Their intrinsic pairing is composition: a character applied to a cocharacter gives \(z\mapsto z^{\langle m,n\rangle}\). The dual torus interchanges these two lattices.
\paragraph{A compatible charge symmetry.}For electric and magnetic labels \((m,n)\in X^*(T)\oplus X_*(T)\), define the alternating pairing
\[\Omega((m,n),(m',n'))=\langle m,n'\rangle-\langle m',n\rangle.\]
It is integral, and the transformation \(S(m,n)=(n,-m)\), interpreted on the dual torus, preserves it: substitution gives \(\langle n,-m'\rangle-\langle n',-m\rangle=\Omega\). Moreover \(S^2=-I\). This verifies the algebraic electric-magnetic exchange at the lattice level without assuming a quantum field theory.
For rank one and a complex parameter \(\tau\) with \(\operatorname{Im}\tau>0\), the positive quadratic expression
\[Q_\tau(m,n)=\frac{|m+\tau n|^2}{\operatorname{Im}\tau}\]
is invariant under \(\tau\mapsto-1/\tau\), \((m,n)\mapsto(n,-m)\): the numerator becomes \(|m+\tau n|^2/|\tau|^2\) and the denominator becomes \(\operatorname{Im}\tau/|\tau|^2\). The positivity assumption excludes division by a vanishing imaginary part.
\paragraph{Missing bridge.}Preservation of a charge pairing and a quadratic form does not construct state spaces, observables, a path integral or a sheaf-theoretic correspondence. In particular no nonabelian group, curve, ramification data or category has been fixed by this benchmark. Establishing that a specified quantum theory realizes the exchange, and that its mathematical objects produce the relevant geometric Langlands equivalence, is the first missing bridge. The broad challenge and unavailable historical source mirror remain as qualified in the supplied statement. The algebra is a standard consistency test, not a new duality theorem.

\paragraph{Verification and adversarial audit.}
Checked pairing signs, S squared, upper-half-plane denominator and distinction between charge arithmetic and categorical duality. This was an LLM adversarial audit, not independent human verification; no theorem prover was used. Prompts A then B applied to this ID: Checked pairing signs, S squared, upper-half-plane denominator and distinction between charge arithmetic and categorical duality. Scope remains limited: Construct a specified quantum theory and mathematical category in which the charge exchange induces the desired correspondence.

Coordinator checked the alternating pairing under the dual-lattice exchange and the rank-one quadratic-form transformation, with positive imaginary part retained. This additional check was model-assisted, not independent human verification.

\paragraph{Tools, process and contributions.}
Python 3 and local shell used to read supplied files and serialize this manuscript; mathematics derived symbolically in the text, with no numerical experiment or formal proof assistant. Contribution prompt A was followed by prompt B. GPT-6 Astra with ultra reasoning drafted the research and performed its stated self-audit; any additional coordinator contribution is identified in the verification and audit. A GPT-6-based Codex coordinator assembled the manuscript and checked repository integration; its exact variant and reasoning configuration were not exposed. No human mathematical contribution, independent human proof check or formal proof certification is supplied. Complete derivations and the consulted source references are included here; private reasoning transcripts are not published.

\paragraph{Final status and first remaining gap.}
\textbf{UNRESOLVED}. Construct a specified quantum theory and mathematical category in which the charge exchange induces the desired correspondence.

\section{Completion Estimate}
\noindent\textbf{COMPLETION: 0\%}
The calculation is an elementary benchmark or conditional reduction; it establishes no claimed new portion of the complete source target.

\subsection{Sources}
\begin{itemize}
\item[\hypertarget{source:statement:1}{S1}] ULAM.AI and the UnsolvedMath Contributors, UnsolvedMath, catalogue entry. Catalogue curation under CC BY 4.0; definitions and English exposition added locally.
\url{https://huggingface.co/datasets/ulamai/UnsolvedMath}.
\item[\hypertarget{source:statement:2}{S2}] DARPA DSO, Mathematical Challenges, BAA 07-68 (catalogue attribution: 2007), Challenge 017. Historical source mirror; unavailable in this check.
\url{https://web.math.utk.edu/~vasili/refs/darpa07.MathChallenges.html}.
\end{itemize}
\label{end:statement}
\end{document}
